\section{Introduction}\label{sec:introduction}

Pricing American-style options requires implied volatility (IV) as an input, yet IV itself is an empirical quantity extracted from observed option prices~\cite{black1973}. This circularity poses a practical problem for synthetic data generation: existing approaches either take IV from the market and feed it into a pricing model, or generate price paths without a consistent IV surface. Parametric IV models such as SABR~\cite{hagan2002} and the SVI parameterization~\cite{gatheral2004} treat the IV surface as an object to be fitted to market quotes, not generated from first principles; they produce smooth interpolations of observed smiles but cannot generate new IV dynamics for unobserved market conditions. Local volatility models~\cite{dupire1994,derman1994} recover a deterministic volatility surface consistent with a single cross-section of option prices, yet the resulting surface is static and cannot produce the stochastic IV fluctuations observed in real markets. Stochastic volatility models such as Heston~\cite{heston1993} generate realistic variance paths, but calibrating them requires observed option prices or IV as input, reintroducing the very circularity that the synthetic generator must avoid. The result is a gap in the literature: no existing framework generates self-consistent IV surfaces as an emergent output of a structural model of equity returns, without requiring observed option data as a calibration target during simulation.

This gap has practical consequences for a growing set of downstream applications. Empirical studies of IV surface dynamics~\cite{cont2002} have established that implied volatility exhibits rich cross-sectional and temporal structure, including smile curvature, skew, and term-structure effects that interact with market regimes. Reproducing this structure in synthetic data is essential for training machine-learning models that must generalize across market conditions, for backtesting hedging strategies such as deep hedging~\cite{buehler2019}, and for conducting scenario-based risk analysis that requires consistent option prices across strikes and expirations. Historical resampling preserves realized distributional features but cannot produce out-of-sample tail scenarios consistent with a structural model of returns; generative approaches based on neural networks have demonstrated promise for equity price paths but do not directly produce self-consistent option price and IV pairs. The scarcity of realistic synthetic option data therefore remains a binding constraint on progress in computational finance.

In this study, we addressed this gap by developing a framework in which IV was an \emph{output} of a structural model of equity returns, not an input. The pipeline coupled a Jump Hidden Markov Model (JumpHMM)~\cite{jumphmm2025} that produced multi-asset price paths with realistic heavy tails, negligible linear autocorrelation, and persistent volatility clustering~\cite{cont2001}, with cross-asset dependence imposed via a Student-$t$ copula~\cite{demarta2005}; a modified Heston stochastic variance process~\cite{heston1993} that converted those price paths into IV paths, where the mean-reversion target $\theta$ was a function of the HMM regime state, days to expiration (DTE), moneyness, and an aggregate market mood indicator; and a Cox-Ross-Rubinstein (CRR) binomial tree~\cite{cox1979} that converted the pre-computed IV into American option prices with early exercise. The central innovation was a hybrid $\theta$-function that linked the Heston mean-reversion target to the JumpHMM regime state, producing regime-dependent volatility dynamics. An equilibrium initialization $v_0 = \theta(t\!=\!0)$ gave each (ticker, strike, DTE) triple its own initial IV, so the smile, skew, and term structure emerged automatically without external IV data. Within this pipeline, the shape function that controlled the smile and term structure of $\theta$ admitted a family of representations, and the modeling choice at each level of the family traded off model capacity for sharing across tickers. We therefore calibrated the shape function through a hierarchy of increasingly expressive forms: a parsimonious five-parameter functional form on single-ticker and cross-ticker semiconductor option chains, a globally shared neural surrogate that replaced the analytic form with a data-driven representation, and a sector-specific neural surrogate that let within-sector shapes specialize. The same decomposition supports a further relaxation to per-ticker networks once per-name sample sizes grow, without changing the downstream pipeline. We demonstrated this progression on a 31-ticker, two-day, six-sector option ladder and showed that the framework produced realistic forward IV paths and American option prices with an endogenous crash-protection premium. The framework is implemented as an open-source Julia package designed to serve as a synthetic option market generator for downstream applications.
